% Part: first-order-logic
% Chapter: natural-deduction
% Section: soundness-identity

\documentclass[../../../include/open-logic-section]{subfiles}

\begin{document}

\olfileid{fol}{ntd}{sid}

\olsection{\usetoken{S}{identity} सह निर्दोषता}

\begin{prop}
$\eq$ साठीच्या नियमांसह नैसर्गिक निगमन निर्दोष आहे.
\end{prop}

\begin{proof}
$\eq[t][t]$ या रूपातील कोणतेही !!{formula} वैध आहे, कारण प्रत्येक
!!{structure}~$\Struct M$ साठी $\Sat{M}{\eq[t][t]}$. (लक्षात घ्या की
$t$ हे बंद पद आहे असे आपण गृहीत धरतो, म्हणजे त्यात कोणतेही चल नाही;
म्हणून चल-विन्यास अप्रासंगिक असतात.)

!!a{derivation} मधील शेवटचे अनुमान \Elim{\eq} आहे असे समजा, म्हणजे
!!{derivation} चे रूप पुढीलप्रमाणे आहे:
\begin{prooftree}
  \AxiomC{$\Gamma_1$}
  \RightLabel{$\delta_1$}
  \DeduceC{$\eq[t_1][t_2]$}
  \AxiomC{$\Gamma_2$}
  \RightLabel{$\delta_2$}
  \DeduceC{$!A(t_1)$}
  \RightLabel{\Elim{\eq}}
  \BinaryInfC{$!A(t_2)$}
\end{prooftree}
आधारविधाने~$\eq[t_1][t_2]$ आणि $!A(t_1)$ ही अनुक्रमे
!!{undischarged} गृहीतके~$\Gamma_1$ आणि $\Gamma_2$ यांच्यापासून
!!{derive} केली आहेत. $!A(t_2)$ हे $\Gamma_1 \cup \Gamma_2$ पासून
निष्पन्न होते, हे दाखवायचे आहे. $\Sat{M}{\Gamma_1 \cup
  \Gamma_2}$ असलेली !!a{structure}~$\Struct{M}$ विचारात घ्या. विगमन
गृहीतकानुसार $\Sat{M}{!A(t_1)}$ आणि $\Sat{M}{\eq[t_1][t_2]}$.
त्यामुळे $\Value{t_1}{M} = \Value{t_2}{M}$. $s$ हा कोणताही चल-विन्यास
आणि $m = \Value{t_1}{M} = \Value{t_2}{M}$ असू द्या.
\olref[fol][syn][ext]{prop:ext-formulas} नुसार
$\Sat{M}{!A(t_1)}[s]$ तेव्हा आणि केवळ तेव्हाच, जेव्हा
$\Sat{M}{!A(x)}[\Subst{s}{m}{x}]$, आणि हे तेव्हा आणि केवळ तेव्हाच, जेव्हा
$\Sat{M}{!A(t_2)}[s]$. $\Sat{M}{!A(t_1)}$ असल्यामुळे
$\Sat{M}{!A(t_2)}$.
\end{proof}

\end{document}
